\psubsection{portability}{\armswitch{0.5}{0.5}{0.2}}{I1}{core,foundry}{One Engine, Two Specifications}
% VI-A. Contribution I1. Evidence: engine report section 1 (sample contract, couplings),
% VF orchestrator (client, patches, in-process path), N5 diff audit, E1/E2 runs, N1 ML runs.
% Wording rules (brief 0.2): never "no code change"; "content digest (SHA-256)", never
% "signed"; the digest covers version and weights only; R09 statistical, R10, R11 and the
% heuristic R07 are not valid measurements; no E8 as the portability result unless withdrawn.
\ifarmwithdrawn
% Withdrawn: P1, P2 and P5 removed; E8 is the only runtime evidence and I1 is downgraded.
Without the tabular arm, I1 holds in a weaker form: the specification is data within one model family, and the tabular specification of \Cref{sec:vera} remains a design.
The runtime evidence is a replay of the language-model campaign under a security-first catalog and registry, covering three requirements on seven deployments with the same runners.
A job served entirely from the call cache reuses the samples of the main campaign, so only the weighting, the recorded digest and the requirement scores differ; \tbd{E8}{foundry}{11 Oct}{number of the 21 E8 jobs served entirely from the cache (cached calls equal to calls, cost 0 EUR)} of the 21 jobs were served that way.
A reweighting of this kind already ran in the earlier study~\apsec.
The replay checks that the property holds on API-served models, which needed a substitute model client and two runner functions patched at run time.
Each run manifest records the content digest (SHA-256) of the catalog, a hash of the registry, the engine revision and the patches in force, so a reviewer can tell which weighting produced a score.
On the language-model panel, \tbd{E1}{foundry}{27 Sep}{LLM headline from valid instruments only (not R09 statistical, R10, R11, heuristic R07), e.g. the requirement with the widest spread across tiers and its range}, and the largest across-seed standard deviation of a requirement score was \tbd{E2}{foundry}{27 Sep}{largest across-seed SD of a requirement score, valid requirements only}; on single-template benchmarks that spread is zero by construction.
\else
% floats/portability.tex -- Tab. 2 (VI-A, contribution I1), column table, 0.15 p.
% Hand-built. Input from 06a-portability.tex whenever the ML arm reports results.
% Established cells, checked on 16 Sep 2026 against the code:
%   digest      catalog_digest() of src/vera/benchmarks/benchmarks_catalog.yaml,
%               recomputed: b25f1485b6df... (also VF config/specs/default_catalog_snapshot.yaml);
%               covers {version, requirement_weights} only (catalog.py:58-70, signing.py:13-20)
%   registry    VF manifests record registry_sha256 (VF orchestrator/vera_adapter.py:245-271)
%   reqs        12 in the catalog; E1 scores 9 per deployment, corpus-scan R03-R05 once
%   instruments 24 distinct ids: 21 benchmarks + 3 dataset scans
%   runners     9 distinct implementations in MVP2_BENCHMARK_REGISTRY, none new
%   outside     VF FoundryLLMClient (VF llm/client.py:1-19); patches lm_eval_limit and
%               ece_logprobs (VF orchestrator/patches.py:1-19), named in every manifest
%   fields      VF calls evaluate_benchmarks in process (VF orchestrator/runner.py:229-270)
% Pending cells: E1 (runs), N1 (ML column), N5 (core diff); registered below the tabular.
\registerfloat{tab:portability}{table}
\begin{table}[t]
  \centering
  \caption{What each specification required of the engine}
  \label{tab:portability}
  \tblsetup
  \begin{tabularx}{\columnwidth}{@{}L{2.45cm}>{\raggedright\arraybackslash}X>{\raggedright\arraybackslash}X@{}}
    \toprule
    & \textbf{Language models} & \textbf{Tabular classifiers} \\
    \midrule
    Catalog digest   & \texttt{b25f14\ldots} & \tbdcell \\
    Requirements     & 12 (9 per model, 3 on a corpus) & \armswitch{4 (\req{03}, \req{05}, \req{07}, \req{11})}{\req{07}, \req{11} run; \req{03}, \req{05} specified}{} \\
    Instruments      & 21 benchmarks, 3 corpus scans & \tbdcell{} criteria, one benchmark each; \tbdcell{} checks \\
    Runners used/new & 9 / 0 & \tbdcell{} / \tbdcell \\
    Code outside engine & client replaced; 2 runners patched at run time & \tbdcell \\
    Core files changed & reference revision & \tbdcell \\
    Fields added     & none & \tbdcell \\
    Runs             & \tbdcell{} of 8 deployments & \armswitch{\tbdcell{} of 3 datasets}{\tbdcell{} on 1 dataset}{} \\
    \bottomrule
  \end{tabularx}
  \par
  {\raggedright
  Requirement identifiers follow COMPL-AI~\cite{guldimann2024complai}. Digest: first six hexadecimal digits of the catalog's SHA-256 content digest (version and weights); run manifests also hash the registry. Core: modules fixed before the tabular arm was coded. Runs: \tbd{E1}{foundry}{27 Sep}{Tab. 2 LLM column: number of the 8 E1 runs completed; confirm from the manifests that the catalog digest is b25f14..., that the recorded VERA revision equals the VF pin and that both patches are listed} and \tbd{N1}{core}{04 Oct}{Tab. 2 ML column: tabular catalog digest, number of fairness criteria and of other checks, runners used and new, runs completed per dataset}. Core diff: \tbd{N5}{core}{05 Oct}{Tab. 2 ML column: core files changed, code outside the engine (predictions producer, launch script), run-context and request fields added; base = revision in the E1 manifests}.\par}
\end{table}
%
% P1: what the ML arm changed in the core (N5).
A reweighting needs only new catalog and registry files~\apsec; a second model family is the stronger test.
Before the tabular arm was coded, we fixed the engine core as the modules that never call a model: catalog loading and its alignment gate, the content digest, the weighted bootstrap over the sample contract of \Cref{sec:vera}, the bands and the triage \decision{D9}{core}{18 Sep}{Which modules form the engine core for the N5 audit: catalog loading and alignment gate, digest, weighted bootstrap, bands, triage; are policy comparison and the sensitivity scripts in or out?}.
Between the revision the language-model runs recorded and the one that ran the tabular arm, \tbd{N5}{core}{05 Oct}{number of core files changed (git diff --stat over the D9 list); I1 may say "no change to the engine core" only if this is zero} core files changed.
The tabular arm added \tbd{N5}{core}{05 Oct}{number of new runners} runners, \tbd{N5}{core}{05 Oct}{number of new registry entries} registry entries and one catalog (\Cref{tab:portability}).

% P2: code outside the core, for both arms (engine report section 1; VF client and patches).
At the language-model revision, the dispatcher was a closed chain of branches that silently sent unknown runner names to a chat-prompt runner.
Neither the run context nor the run request had a slot for labels, predictions or a protected attribute.
The tabular arm \tbd{N5}{core}{05 Oct}{how each coupling was passed: dispatcher entries added and whether unknown names now raise (G2); run-context and request fields added (G1)}.
The language-model arm needed code too: the engine's client sends every call to a configured Ollama endpoint, so the Foundry campaign substitutes a client with the same signature and patches two runner functions at run time.
One patch lifts a ten-item cap on evaluation-harness tasks; the other reads calibration confidence from token log-probabilities where a deployment returns them, and otherwise flags the heuristic confidence as a fallback.

% P3: same scoring path, digests and provenance (VF runner.py, vera_adapter.py).
The Foundry campaign scores in process with the engine's dispatcher, weighted bootstrap and band settings.
Each run manifest records the content digest (SHA-256) of the catalog, a hash of the registry, the engine revision and the patches in force.
The catalog digest covers weights but not band thresholds \decision{D29}{core}{25 Sep}{G5: add per-criterion thresholds and the registry to the catalog and its digest, or describe the environment-set thresholds as they are?}.
The tabular runs \tbd{N1}{core}{04 Oct}{entry point of the ML runs (script or campaign CLI), whether they call the same bootstrap and bands, and which digests and provenance fields their manifests record}.

% P4: arm L results (E1, E2).
On the language-model panel, \tbd{E1}{foundry}{27 Sep}{LLM headline from valid instruments only (not R09 statistical, R10, R11, heuristic R07), e.g. the requirement with the widest spread across the three tiers and its range}.
Across three seeds on four deployments, the largest standard deviation of a requirement score was \tbd{E2}{foundry}{27 Sep}{largest across-seed SD of a requirement score, valid requirements only}.
On single-template benchmarks this spread is zero by construction.

% P5: arm M results (N1), by scenario.
\ifarmreduced
The reduced arm ran fairness and calibration on one dataset, \tbd{N1}{core}{04 Oct}{name of the single dataset the reduced arm ran on}; data adequacy and privacy are specified but were not run.
On that dataset, \tbd{N1}{core}{04 Oct}{reduced-arm result: which fairness criteria ran, whether the engine flagged the disparity (recall against the manifest if synthetic), continuous values}.
I1 therefore rests on a language-model panel and one tabular demonstration.
\else
\parplan[20]{On the synthetic credit corpus, the engine's findings are checked against the manifest of the planted disparity and proxy; the public datasets are scored with the same catalog.}{N1 runs due 4 Oct (G3, G9). Report recall of the planted disparate impact, and of the planted proxy only if D3 keeps G3b (22 Sep), against the generator manifest. Then per-criterion values on Adult (becker1996adult) and South German Credit (gromping2019south with hofmann1994statlog; D6), as continuous values, never thresholded 0/1. Detection and scoring only: which verdicts are criterion-robust belongs to VI-B.}
\fi
\fi
